\begin{figure*}[t]
\centering
\setlength{\tabcolsep}{0.8pt}
\renewcommand{\arraystretch}{0.5}
\newcommand{\cw}{0.131\textwidth}
\begin{tabular}{@{}ccccccc@{}}
\small テストビュー & \small 正解画像 & \small \ours（本手法） & \small SSD-GS & \small \gsthree & \small OLAT-GS & \small RNG \\[2pt]
\includegraphics[width=\cw]{figures/comparison/translucent/full.png} & \includegraphics[width=\cw]{figures/comparison/translucent/GT.png} & \includegraphics[width=\cw]{figures/comparison/translucent/LiSA-staged.png} & \includegraphics[width=\cw]{figures/comparison/translucent/SSD-GS.png} & \includegraphics[width=\cw]{figures/comparison/translucent/GS3.png} & \includegraphics[width=\cw]{figures/comparison/translucent/OLAT-Gaussians.png} & \includegraphics[width=\cw]{figures/comparison/translucent/RNG.png} \\
\footnotesize Translucent & \footnotesize PSNR & \footnotesize \textbf{34.36} & \footnotesize 32.28 & \footnotesize 32.58 & \footnotesize 29.93 & \footnotesize 27.92 \\[3pt]
\includegraphics[width=\cw]{figures/comparison/hotdog/full.png} & \includegraphics[width=\cw]{figures/comparison/hotdog/GT.png} & \includegraphics[width=\cw]{figures/comparison/hotdog/LiSA-staged.png} & \includegraphics[width=\cw]{figures/comparison/hotdog/SSD-GS.png} & \includegraphics[width=\cw]{figures/comparison/hotdog/GS3.png} & \includegraphics[width=\cw]{figures/comparison/hotdog/OLAT-Gaussians.png} & \includegraphics[width=\cw]{figures/comparison/hotdog/RNG.png} \\
\footnotesize Hotdog & \footnotesize PSNR & \footnotesize \textbf{33.97} & \footnotesize 31.93 & \footnotesize 32.08 & \footnotesize 26.85 & \footnotesize 28.49 \\[3pt]
\includegraphics[width=\cw]{figures/comparison/pikachu/full.png} & \includegraphics[width=\cw]{figures/comparison/pikachu/GT.png} & \includegraphics[width=\cw]{figures/comparison/pikachu/LiSA-staged.png} & \includegraphics[width=\cw]{figures/comparison/pikachu/SSD-GS.png} & \includegraphics[width=\cw]{figures/comparison/pikachu/GS3.png} & \includegraphics[width=\cw]{figures/comparison/pikachu/OLAT-Gaussians.png} & \includegraphics[width=\cw]{figures/comparison/pikachu/RNG.png} \\
\footnotesize Pikachu（実撮影） & \footnotesize PSNR & \footnotesize 30.44 & \footnotesize 29.58 & \footnotesize \textbf{30.71} & \footnotesize 29.46 & \footnotesize 16.31 \\[3pt]
\includegraphics[width=\cw]{figures/comparison/fail_drums/full.png} & \includegraphics[width=\cw]{figures/comparison/fail_drums/GT.png} & \includegraphics[width=\cw]{figures/comparison/fail_drums/LiSA-staged.png} & \includegraphics[width=\cw]{figures/comparison/fail_drums/SSD-GS.png} & \includegraphics[width=\cw]{figures/comparison/fail_drums/GS3.png} & \includegraphics[width=\cw]{figures/comparison/fail_drums/OLAT-Gaussians.png} & \includegraphics[width=\cw]{figures/comparison/fail_drums/RNG.png} \\
\footnotesize Drums & \footnotesize PSNR & \footnotesize 30.90 & \footnotesize 31.32 & \footnotesize \textbf{31.93} & \footnotesize 22.23 & \footnotesize 16.23 \\[3pt]
\end{tabular}
\caption{\textbf{定性的比較。}テストビューの枠内を切り出して示す（PSNR はビュー全体）。各画像は、そのシーンで最良のベースラインに対する \ours のビューごとの優位性が中央値となるビューである。上から、投影される影が支配的な二つの \gsthree シーン、各手法の差が小さい実撮影シーン、本手法が最も劣る合成シーンを示す。SSD-GS と \gsthree のプリミティブごとの影は、花瓶の取っ手とホットドッグの移動する影を断片化するが、アトラスを参照する画素受光点はこれらを再現する。一方、\ours は Drums のシンバルの細い光の筋を捉えられない。Fish を含む他のシーンは補足資料に示す。}
\label{fig:comparison}
\end{figure*}
